\section{Introduction}\label{sec:introduction}

The Birch--Swinnerton-Dyer conjecture predicts that the order of vanishing
of the $L$-function of an elliptic curve equals its Mordell--Weil rank,
and that its Shafarevich--Tate group is finite; its refined form also
specifies the leading coefficient \cite{BirchSwinnertonDyer1965,Tate1966}.
For analytic order zero or one, the rank equality and finiteness are
theorems of Gross--Zagier and Kolyvagin
\cite{GrossZagier1986,Kolyvagin1990}.
The converse problem begins with algebraic information and seeks the
corresponding analytic order.

The appropriate algebraic hypothesis here is the corank of the full
$2$-power Selmer group. It records a possible divisible contribution
from the Shafarevich--Tate group as well as the Mordell--Weil rank.
With the usual local Kummer conditions, put
\[
s_2(E)=\corank_{\Z_2}\Sel_{2^\infty}(E/\Q).
\]
Our main result is the following pointwise statement.

\begin{theorem}[The rational-two-torsion $2$-converse]\label{thm:main}
Let $E/\Q$ be an elliptic curve with $E(\Q)[2]\ne0$ and
$s_2(E)\in\{0,1\}$. Then
\[
\ord_{s=1}L(E,s)=\rank_{\Z}E(\Q)=s_2(E),
\qquad
\Sha(E/\Q)\text{ is finite}.
\]
\end{theorem}

The statement includes one rational two-torsion line and full rational
two-torsion, CM and non-CM curves, every rational isogeny configuration,
and every reduction type at $2$ and at the bad primes. It is a
$2$-converse in this residual range. It does not assert the full
Birch--Swinnerton-Dyer leading-coefficient formula, nor does a bound
$\rank E(\Q)\le1$ alone imply its Selmer hypothesis.

The broader companion \cite[Theorem~1.1]{OpenAIGoldfeldTwoConverse2026}
proves the $2$-converse without the rational-two-torsion hypothesis.
Here we give a separate full treatment of the rational-two-torsion case,
developing its integral interpolation and finite twist constructions
within this paper; the broader converse is not an input to the proof.

The converse with the stronger hypotheses
$\rank E(\Q)=1$ and $\Sha(E/\Q)$ finite is known for every
elliptic curve over $\Q$; see \cite[\S1]{Kim2023SoftConverse}.
The fixed-prime problem asks instead whether a low Selmer corank
alone forces the analytic order. Skinner and Wei Zhang established
general rank-one converses at good ordinary odd primes under residual
and local hypotheses \cite{Skinner2020,Zhang2014Heegner}.
For odd Eisenstein primes, Castella, Grossi, Lee and Skinner prove
the good-reduction converse under a local-character restriction
\cite[Theorem~E]{CGLS2022}; Keller and Yin treat potentially good
ordinary reduction without that restriction
\cite[Theorem~B]{KellerYin2024}. Smith's theorem gives the $50/50$
distribution of $2^\infty$-Selmer coranks in every fixed quadratic-twist
family, counting twist parameters of both signs \cite{Smith2025};
this density statement does not decide a prescribed member.

The CM case includes pointwise results at the prime $2$.
Burungale and Tian prove the rank-zero $p$-converse for CM curves at
every prime \cite{BurungaleTian2026}. In rank one, Burungale and
Skinner treat good ordinary primes and, with a localization hypothesis,
good supersingular primes
\cite[Theorems~A.1 and~A.3, Remark~A.5]{BurungaleSkinner2024Appendix}.
The earlier ordinary CM converse of Burungale, Castella, Skinner and
Tian retains a Hecke-conductor condition
\cite{BurungaleCastellaSkinnerTian2022}.
Kriz's preprint gives rank-zero and rank-one converses for maximal-order
CM curves at primes ramified in the CM field, including the pointwise
$2$-converse for the congruent-number family
\cite[Theorems~1.1 and~10.17(2)]{Kriz2022CM}.
These results already treat individual curves in their stated ranges;
the theorem here includes non-CM curves with rational two-torsion and
imposes no reduction condition at $2$.

The proof brings together two arithmetic interpolation constructions and a
finite construction with quadratic twists.  Kato's zeta classes provide
the cyclotomic input \cite{Kato2004}; Ferrero--Washington's
vanishing theorem controls the residual cyclotomic cohomology
\cite{FerreroWashington1979}.  Heegner points and their derivative
classes provide the ring-class input, following the Euler-system method
of Kolyvagin and its cohomological formulation by Howard
\cite{Kolyvagin1990,Howard2012}.  We retain integral local conditions and
duality maps in the framework of Selmer complexes \cite{Nekovar2006}.
The coefficient tests use the half-integral weight theory of Shimura
and Waldspurger's relation between coefficients and central values
\cite{Shimura1973HalfIntegral,Waldspurger1981}.  Their ternary theta
construction lies in the definite quaternion framework developed by
Gross \cite{Gross1987Heights}; we use the Weil representation and the
reduction description of Jetchev and Kane for the weighted Heegner-point
tests \cite{Weil1964,JetchevKane2010}.
Each of these inputs has a precise role; none supplies the pointwise
conclusion by itself.

Here is the common interpolation problem.  Fix the curve that is to be
tested, and index a finite family of quadratic twists by
$x\in\mathbb F_2^b$, with $x=0$ denoting that curve.  Each auxiliary
prime divides the twist parameter according to a linear form in $x$.
We construct these families so that every nonzero address has the
required analytic order and a uniform normalized valuation bound.
An integral interpolation argument then recovers the value or derivative
at the missing address $0$.  The number and sizes of the auxiliary
primes may grow with $b$; the valuation bound may not.

For Selmer corank zero, the required nonzero values are central
$L$-values.  Their normalized valuations are measured by coefficients
of half-integral weight forms.  The cyclotomic determinant makes these
values part of one integral family, and a congruence on a sufficiently
large binary cube recovers the missing value.  This also supplies the
even estimate needed to normalize the rank-one construction.

The rank-one construction requires a further preparation.  Fix a negative
quadratic discriminant $h_*$ with $L(E^{(h_*)},1)\ne0$.
For positive $h$ in the prescribed local squareclasses, the
genus-character Heegner sum $P_h$ attached to $\Q(\sqrt{hh_*})$
detects a simple zero of $E^{(h)}$: its height measures
$L'(E^{(h)},1)L(E^{(h_*)},1)$.  To normalize the corresponding coefficient tests
uniformly, we first bound the $2$-divisibility of $P_h$ from below
in terms of the prime factors of $h$, with a constant independent
of their number.
For each parameter $h$ of known analytic rank one, a different negative companion
$k'(h)$, with a bounded number of prime factors, combines the even
estimate with a ring-class lower bound to give this uniform bound.
Only after that argument do we normalize the odd coefficients and use
their minimum-depth values as binary nonvanishing tests.

For the return to a prescribed corank-one curve, choose one final
negative discriminant $k$ with $L(E^{(k)},1)\ne0$, and fix
$K=\Q(\sqrt{k})$.  This companion does not vary with $h$ or with the
cube dimension.  The quadratic corank identity gives $s_2(E/K)=1$.
We construct families with every new prime split in $K$ for which,
at every nonzero address, $E^{(h)}$ has a simple zero and $E^{(hk)}$
has nonzero central value.  The two coefficient tests together bound
the normalized index of the ring-class Heegner sum over $K$.
Ring-class interpolation then recovers the derivative at $h=1$.
Thus $h_*$ defines the detector, $k'(h)$ proves its lower bound,
and the fixed $k$ permits the final transfer.

The arithmetic difficulty is uniform integral control at the prime $2$.
Residual irreducibility is unavailable here.  On the cyclotomic side we
use Kato classes from one fixed modular-form lattice, separating rational
divisibility from uniform integrality.  On the ring-class side we retain
the full integral Kummer conditions at the bad places and actual maps
from finite free models to Galois cochains.  Heegner derivative operations
preserve the determinant valuation even when a localized class is
divisible by $2$.  These constructions give the missing-vertex
Theorems~\ref{cy:missing-vertex} and~\ref{rd:missing-vertex}.

The finite construction has a different task: it must make the required
coefficient tests nonzero at every nonzero address simultaneously.
The variables are quadratic residue symbols between auxiliary primes.
Changing one such symbol contracts two prime labels; repeated changes
are governed by forests.  Finite networks of new primes then isolate
one address at a time.  For rank one, the derivative test and its
central-value partner may be separately nonzero without taking nonzero
values on the same assignment of residue symbols.  Theorem~\ref{gr:address}
adjoins prime networks making both tests equal to $1$ at every nonzero
address.

Figure~\ref{fig:proof-roadmap} records the order in which the valuation
bounds are obtained and used.  Section~\ref{sec:preliminaries} fixes the
Selmer and local conventions.  Section~\ref{cy:section} proves the
cyclotomic interpolation statements; Sections~\ref{rl:section}
and~\ref{rd:section} construct the ring-class deformation and prove its
bounds and transfer theorem.  Section~\ref{sec:nonvanishing} supplies
auxiliary twists with prescribed local behavior and a bounded number of
prime factors where needed.  Sections~\ref{co:section}
and~\ref{gr:section} give the coefficient tests and finite networks.
Section~\ref{pw:section} assembles these results at the prescribed curve.

\begin{figure}[htbp]
\centering
\small
\setlength{\tabcolsep}{6pt}
\begin{tabular}{c}
\fbox{\parbox{0.86\textwidth}{\centering
Even coefficients, finite prime networks and cyclotomic transfer\\
\begin{tabular}{@{}c@{\qquad}c@{}}
$\swarrow$ & $\searrow$\\
Uniform even estimate & Corank-zero conclusion\\
(Proposition~\ref{pw:partitionable-even}) & (Proposition~\ref{pw:even-base})
\end{tabular}}}
\\[5pt]
Even estimate $\Downarrow$\\[5pt]
\fbox{\parbox{0.86\textwidth}{\centering
Fix $h_*$ to define $P_h$; choose a bounded-factor companion $k'(h)$\\
Even estimate + ring-class lower bound\\
$\Downarrow$\\
Uniform lower bound for the $2$-divisibility of $P_h$\\
and odd coefficient normalization}}
\\[5pt]
$\Downarrow$\\[5pt]
\fbox{\parbox{0.86\textwidth}{\centering
Fix $k$ with $L(E^{(k)},1)\ne0$ and put $K=\Q(\sqrt{k})$\\
Odd and even tests + finite prime networks split in $K$\\
$\Downarrow$\\
Simple zeros for $E^{(h_x)}$ and nonzero values for $E^{(h_xk)}$\\
with a common normalized index bound, for every $x\ne0$}}
\\[5pt]
$\Downarrow$\\[5pt]
\fbox{\parbox{0.86\textwidth}{\centering
Fixed-$K$ ring-class transfer and $s_2(E/K)=1$\\
$\Downarrow$\\
Analytic rank one at the original vertex $x=0$}}
\end{tabular}
\caption{The even construction proves the corank-zero case and a uniform
estimate; the latter supplies the bound needed to normalize the odd coefficients.
The detector partner $h_*$ is fixed, the bounded-factor companion $k'(h)$
varies only in the lower-bound step, and the final companion $k$ is fixed
before the simultaneous families are constructed.  The families and
auxiliary primes may grow with the cube dimension; the normalized
index bound may not.}
\label{fig:proof-roadmap}
\end{figure}

